\begin{abcliste}
\abc
Das Drehmoment ist $M = F\cdot d$, mit dem Hebelarm $d$: dem Abstand von D zur \emph{Wirkungslinie} der Kraft, nicht zum Angriffspunkt.
\begin{itemize}
 \item $F_1$ wirkt senkrecht; ihre Wirkungslinie ist die senkrechte Gerade 3 Kästchen rechts von D, also $d_1 = \daO$.
 \item $F_2$ wirkt waagrecht; ihre Wirkungslinie ist die waagrechte Gerade 1 Kästchen über D, also $d_2 = \dbO$. Dass sie 4 Kästchen rechts von D angreift, spielt keine Rolle: Man kann die Kraft entlang ihrer Wirkungslinie verschieben.
 \item $F_3$ wirkt senkrecht, 5 Kästchen links von D: $d_3 = \dcO$.
\end{itemize}
\begin{align}
 M_1 &= \MaF = \Fa\cdot\da = \Ma \approx \result{\MaP{0}{2}} \quad\text{(Gegenuhrzeigersinn)}\\
 M_2 &= \MbF = \Fb\cdot\db = \Mb \approx \result{\MbP{0}{2}} \quad\text{(Uhrzeigersinn)}\\
 M_3 &= \McF = \Fc\cdot\dc = \Mc \approx \result{\McP{0}{2}} \quad\text{(Gegenuhrzeigersinn)}
\end{align}
$F_1$ drückt rechts von D nach oben und $F_3$ links von D nach unten: Beide drehen die Platte im Gegenuhrzeigersinn. $F_2$ zieht oberhalb von D nach rechts und dreht sie im Uhrzeigersinn. Die Reihenfolge ist
\[
 \result{M_2 < M_3 < M_1}
\]

\abc
Lea ordnet nach den Kräften: $F_3 < F_1 < F_2$. Ein Drehmoment hängt aber nicht nur von der Kraft, sondern auch vom Hebelarm ab: Die grösste Kraft $F_2$ hat den kürzesten Hebelarm und das kleinste Drehmoment.

Tim nimmt bei $F_2$ den Abstand von D zum Angriffspunkt:
\begin{align}
 r_2 &= \rbF = \sqrt{\left(\xb\right)^2+\left(\db\right)^2} = \rb \\
 F_2\cdot r_2 &= \Fb\cdot\rb = \Mw \approx \MwP{0}{2}
\end{align}
Er erhält ebenfalls $F_3 < F_1 < F_2$ (bei $F_1$ und $F_3$ liegt der Angriffspunkt auf gleicher Höhe wie D; dort ist der Abstand zum Angriffspunkt zufällig gleich dem Hebelarm). Für $F_2$ ist $r_2$ aber nicht der Hebelarm: Der Hebelarm ist das Lot von D auf die Wirkungslinie, hier nur \dbO.
\end{abcliste}
